\begin{tablalafrox}{Diccionario del Producto 8: Infraestructura, equipamiento y servicios habilitados. Fuente: elaboración propia de los paquetes de infraestructura y equipamiento a partir del diseño físico de la solución propuesta.}{tab:T14-dic-prod-8}
  {L{1.2cm} Y Y Y C{0.9cm} L{2.3cm}}
  {\cab{Código} & \cab{Paquete de trabajo} & \cab{Entregable} & \cab{Criterio de aceptación} & \cab{Resp.} & \cab{Período}}
  \multicolumn{6}{>{\columncolor{lafroxBg}\RaggedRight\arraybackslash}p{\dimexpr\textwidth-2\tabcolsep\relax}}{\textbf{Cuenta 8.1 — Compras de hardware y especificaciones.} Esta cuenta de control reúne la especificación del equipamiento de terreno que compra el CLIENTE, la especificación y compra de la infraestructura que provee LafroX y su recepción técnica. LafroX provee, instala, integra y mantiene la infraestructura dentro del precio contratado (BA, art. 14.2, p. 10); su valorización corresponde a la Oferta Económica conforme al artículo 51 de las mismas Bases (p. 29).} \tabularnewline*
  8.1.1 & Especificación de compra de equipos de uso para el CLIENTE & Especificación de terminales, impresoras, terminales de pago y balanzas. & Modelos, cantidades y necesidad por ola aprobados por el CLIENTE; el sistema de frío corresponde a 8.1.4. & ARQ & E1 · Mes 2. \\
  8.1.2 & Especificación y compra de la sala técnica, los racks, los servidores y los gabinetes de borde & Especificación y orden de compra emitida por LafroX. & Coincide con los planos de la cuenta 2.4 y con el T-11. LafroX emite la orden de compra en el mes 2 para instalar la sala en el mes 3. Incluye la red inalámbrica del T-11; la reserva cubre la diferencia que fije 2.2.2. & ARQ & E1 · Mes 2. \\
  8.1.3 & Actas de recepción técnica de equipamiento de terreno e infraestructura & Actas de recepción técnica del equipamiento de terreno comprado por el CLIENTE y de la infraestructura provista por LafroX. & El 100 \% de lo especificado se recibe conforme antes de su instalación en los sitios (cuentas 8.3 a 8.7). & SRE & E1 · Según el calendario de compra. \\
  8.1.4 & Especificación y compra del sistema de cadena de frío & Especificación de compra para el CLIENTE (B-01 y kits de camión) y orden de compra de LafroX para los gateways B-02 y el prototipo. & B-01/B-02 de cámaras y componentes para 31 kits de camión; cotización local y certificación Subtel aplicable antes del banco; compra masiva del CLIENTE sujeta a aceptación del prototipo. & ARQ & E1 · Meses 1–2; antes del banco. \\
  \midrule
  \multicolumn{6}{>{\columncolor{lafroxBg}\RaggedRight\arraybackslash}p{\dimexpr\textwidth-2\tabcolsep\relax}}{\textbf{Cuenta 8.2 — Servicios de nube, licencias y enlaces contratados.} Esta cuenta reúne los contratos de servicios de nube, licencias y enlaces de los sitios que suscribe LafroX. El caso identifica la dependencia de un solo enlace en Concepción como un riesgo, y Los Ángeles tiene señal intermitente en su ventana de madrugada. Cada CD tiene fibra como enlace principal, LTE de respaldo y Starlink como tercer camino en espera caliente. Cada plataforma de cross-docking tiene Starlink como enlace principal y LTE de dos proveedores como respaldo.} \tabularnewline*
  8.2.1 & Contrato de los servicios de AWS & Servicios contratados y cuentas habilitadas. & Todos los servicios de la arquitectura física del Capítulo 4 están activos antes de configurar los ambientes (3.1). & SRE & E1 · Antes del mes 4. \\
  8.2.2 & Suscripción del servicio de gestión de dispositivos y de detección en endpoints & Suscripciones activas. & Las licencias cubren todos los equipos del T-11 antes de su enrolamiento. & SRE & E1 · Antes del enrolamiento de los terminales. \\
  8.2.3 & Licencias de software de terceros constituidas a nombre del CLIENTE & Licencias a nombre del CLIENTE y su registro. & Ninguna licencia de terceros está a nombre de LafroX, y todas están vigentes. & ARQ & E1 · Antes de usar cada producto. \\
  8.2.4 & Contrato de fibra óptica de los CD de Talca y Concepción & Contrato y fecha de instalación. & La capacidad contratada coincide con la del T-11. & SRE & E1 · Antes del H3. \\
  8.2.5 & Contrato de planes LTE con dos proveedores (ocho planes) & Contratos. & Los ocho planes cubren un respaldo LTE en Talca, uno en Concepción y dos de proveedores distintos en cada plataforma de cross-docking. & SRE & E1 · Antes del H3. \\
  8.2.6 & Contrato de Starlink para los cinco sitios & Contratos y equipos de Starlink para los cinco sitios (Talca, Concepción y las tres plataformas de cross-docking). & Cubren los 56 meses del contrato en los cinco sitios. & SRE & E1 · Antes del H3. \\
  8.2.7 & Contrato de SIM M2M y servicio satelital de los camiones con frío & Contratos de dos operadores móviles y Globalsat Chile para Iridium SBD & 62 SIM M2M, servicio bidireccional y camino de ingesta acordado antes del piloto; sin confundir con SIM de terminales ni enlaces de sitio. & SRE & E1 · Mes 2; servicio disponible antes del banco y del piloto. \\
  \midrule
  \multicolumn{6}{>{\columncolor{lafroxBg}\RaggedRight\arraybackslash}p{\dimexpr\textwidth-2\tabcolsep\relax}}{\textbf{Cuenta 8.3 — Sala técnica de Talca habilitada.} Esta cuenta de control reúne el trabajo de convertir la sala de Talca en una sala técnica que cumpla las Bases (BTT, cap. 6, pp. 14--17). Termina con un acta que habilita el montaje de los racks.} \tabularnewline*
  8.3.1 & Piso técnico de la sala y recepción de obra civil de separación (BTT, cap. 6, RT-06.06, p. 15) & Sala acondicionada según el plano. La obra civil de separación está a cargo del CLIENTE (BTT, cap. 6, RT-06.06, p. 15); LafroX la especifica y coordina, y provee el piso técnico. & El recinto coincide con el plano aprobado, incluido el blindaje perimetral especificado, y está listo para recibir los equipos. & SRE & E1 · Mes 3, después de la especificación 8.1.2 (D-10). \\
  8.3.2 & Sistema de energía de la sala instalado (UPS N+1, generador y PDU) & Sistema de energía instalado y probado: UPS modular N+1, generador con 24 horas de autonomía, transferencia automática y PDU A/B. & En una prueba con carga real, el corte de la red pasa a la UPS y al generador sin que se apague ningún equipo. & SRE & E1 · Mes 3. \\
  8.3.3 & Hardware de climatización de precisión N+1 instalado & Climatización instalada y probada. & Con un equipo apagado, la sala se mantiene en el rango de temperatura del fabricante. & SRE & E1 · Mes 3. \\
  8.3.4 & Hardware de detección temprana de incendio y extinción por agente limpio instalado & Sistema instalado y conectado al monitoreo. & Una prueba de detección genera la alarma en el monitoreo, y la extinción queda armada. & SRE & E1 · Mes 3. \\
  8.3.5 & Acta de recepción técnica de la sala & Acta firmada con disponibilidad del recinto. & Obra civil y piso técnico 8.3.1 recibidos y pruebas 8.3.2 a 8.3.4 aprobadas. La disponibilidad del espacio por el CLIENTE se acredita mediante el acta firmada del recinto antes del montaje. & SRE & E1 · Mes 4, después de 8.3.1 a 8.3.4 y antes del montaje de los racks (8.4). \\
  \midrule
  \multicolumn{6}{>{\columncolor{lafroxBg}\RaggedRight\arraybackslash}p{\dimexpr\textwidth-2\tabcolsep\relax}}{\textbf{Cuenta 8.4 — Cableado, racks y gabinetes.} Esta cuenta reúne la red física de cables de la sala y las bodegas y el montaje de los equipos de cómputo y comunicaciones en cada uno de los cinco sitios.} \tabularnewline*
  8.4.1 & Cableado de cobre Cat6A y fibra OM4 tendido, con dos ductos de ingreso independientes & Cableado tendido. & Los dos ductos de ingreso están separados (BTT, cap. 6, RT-06.32, p. 17), y todos los puntos del diseño están cableados. & SRE & E1 · Mes 4. \\
  8.4.2 & Certificación de enlaces, etiquetado y planos as-built & Informe de certificación por enlace y planos as-built. & El 100 \% de los enlaces está certificado y etiquetado. & SRE & E1 · Mes 6. \\
  8.4.3 & Gabinetes de piso y switches de acceso instalados en las bodegas de Talca y Concepción & Gabinetes de piso y puntos de acceso instalados y conectados. & Cada punto tiene red y la cobertura medida cumple el SD4, Anexo 4-W.13. & SRE & E1 · Antes del H6. \\
  8.4.4 & Rack R01 de servidores montado y traslado del servidor ERP & Rack R01 montado (clúster de tres nodos, respaldo local y consola KVM), energizado por los dos caminos y conectado. Servidor del ERP de 2017, propiedad del CLIENTE, trasladado desde la sala actual de 25 m² a R01 (U25–U26), con respaldo completo verificado, en ventana dominical y sin cambios de software. & Cada equipo tiene doble fuente en circuitos distintos. El clúster sigue funcionando con un nodo apagado. El traslado del ERP a R01 se realiza después del H3 y antes de la marcha blanca de Talca, con respaldo completo verificado, en ventana dominical y sin cambios de software (SD4, 4.3.1.4; T-11). & SRE & E1 · Montaje aceptable desde el 9 de junio de 2027. Traslado en la primera ventana dominical admisible posterior al acta H3, con registro posterior y antes de la marcha blanca de Talca. \\
  8.4.5 & Rack R02 de comunicaciones montado & Rack R02 montado y conectado (firewalls, switches de núcleo y de gestión, y distribución de fibra). & Al apagar un firewall o un switch, el tráfico sigue sin cortarse. & SRE & E1 · Antes del H3. \\
  8.4.6 & Gabinete de borde de Concepción montado & Gabinete montado y monitoreado, con dos servidores HPE DL20 (activo y en espera), firewalls, switches, UPS, climatización y monitoreo. & El gabinete reporta su estado al centro de operaciones, la falla de un firewall no corta la red y el servidor de espera toma el servicio si falla el activo. & SRE & E1 · Antes del H3. \\
  8.4.7 & Gabinetes de borde de plataformas de cross-docking montados & Tres gabinetes montados en Curicó, Chillán y Los Ángeles, cada uno con dos mini-PC industriales E-01 (activo y en espera). & Cada plataforma opera su ventana de madrugada con el enlace cortado y el mini-PC de espera toma el servicio si falla el activo. La cobertura medida cumple el SD4, Anexo 4-W.13. & SRE & E1 · Antes de la ola de cross-docking. \\
  \midrule
  \multicolumn{6}{>{\columncolor{lafroxBg}\RaggedRight\arraybackslash}p{\dimexpr\textwidth-2\tabcolsep\relax}}{\textbf{Cuenta 8.5 — Seguridad física.} Esta cuenta de control reúne lo que protege físicamente a la sala y a sus respaldos.} \tabularnewline*
  8.5.1 & Hardware de control de acceso biométrico y esclusa instalado & Control de acceso instalado y con las personas autorizadas enroladas. & Solo las personas enroladas entran, y cada ingreso queda registrado. & SEG & E1 · Antes del H3. \\
  8.5.2 & Hardware de videovigilancia instalado & Cámaras instaladas y conectadas al monitoreo. & Se pueden consultar las imágenes de las seis zonas de los últimos 30 días. & SEG & E1 · Antes del H3. \\
  8.5.3 & Recinto de custodia de respaldos habilitado & Recinto habilitado. & El recinto mantiene entre 18 y 27 °C y entre 40 y 60 \% de humedad. & SEG & E1 · Antes del H3. \\
  \midrule
  \multicolumn{6}{>{\columncolor{lafroxBg}\RaggedRight\arraybackslash}p{\dimexpr\textwidth-2\tabcolsep\relax}}{\textbf{Cuenta 8.6 — Equipamiento de campo.} Esta cuenta reúne la puesta en marcha, por olas, de los equipos utilizados en bodega y ruta, con instalación previa al comienzo de cada ola.} \tabularnewline*
  8.6.1 & Terminales de preventa, reparto y bodega configurados y enrolados & Terminales enrolados y asignados. & El 100 \% de los terminales de cada ola está enrolado y probado antes de que la ola empiece. & SRE & E1 · Antes de cada ola (H6). \\
  8.6.2 & Terminales de pago configurados, uno por camión & Terminales de pago configurados y asignados a cada camión. & Un cobro de prueba con tarjeta se autoriza y queda unido a su entrega. & SRE & E1 · Antes de la ola de reparto. \\
  8.6.3 & Sensores de cámara y gateways IoT instalados en Talca y Concepción & Sensores y gateways instalados y calibrados. & Los 21 puntos reportan a M12 sin brechas durante una semana antes de la marcha blanca. & SRE & E1 · Antes del H5. \\
  8.6.4 & Kit de cadena de frío instalado y recibido en los 28 camiones con frío & Acta por camión, kit instalado y chequeo completo. & Cuatro actas antes del piloto y 28 antes de la ola 3; alarma simulada, asociación de registradores y protección local comprobadas. & SRE & E1 · Meses 8–12; antes de la ola de reparto. \\
  8.6.5 & Impresoras de cabina instaladas & Impresoras instaladas y vinculadas al terminal de reparto. & Cada camión imprime sin señal. & SRE & E1 · Antes de la ola de reparto. \\
  8.6.6 & Impresoras de andén y balanzas de recepción instaladas & Equipos instalados y conectados. & La recepción imprime etiquetas y registra el peso desde M1. & SRE & E1 · Antes de la ola de recepción. \\
  \midrule
  \multicolumn{6}{>{\columncolor{lafroxBg}\RaggedRight\arraybackslash}p{\dimexpr\textwidth-2\tabcolsep\relax}}{\textbf{Cuenta 8.7 — Configuración de los sitios.} Esta cuenta de control reúne lo necesario para que el hardware instalado funcione como plataforma: el software de base, los enlaces y la prueba de autonomía.} \tabularnewline*
  8.7.1 & Software de base instalado en los cinco sitios & Software de base instalado en los cinco sitios (virtualización, contenedores, PostgreSQL, RabbitMQ y Keycloak). & Cada sitio ejecuta la misma imagen de la aplicación que la nube. & SRE & E1 · Antes del H3. \\
  8.7.2 & Enlaces de fibra, LTE y Starlink habilitados, con conmutación al respaldo & Enlaces operativos. & Al cortar el enlace principal, el respaldo toma el tráfico en menos de 30 segundos. & SRE & E1 · Antes del H3. \\
  8.7.3 & Borde de los CD en servicio, con la prueba de autonomía de 24 horas aprobada & Informe de la prueba de autonomía por sitio. & Cada CD opera 24 horas sin enlace y concilia sin pérdidas ni duplicados. & SRE & E1 · Evidencia de autonomía al cierre del 18 de junio de 2027, previa al ingreso del expediente H3; pruebas ampliadas antes del H5. \\
\end{tablalafrox}
